\section{Keeping the unwanted primes out}\label{sec:spectral}
A norm representation is only a candidate. A prime $p\equiv2\pmod3$
can occur to even exponent in an Eisenstein norm, and that already keeps
the value out of $\AA$. The useful feature is that divisibility of the
norm forces divisibility of both coordinates.

For $p\in\PP$, the quotient $R/pR$ is a field, including when $p=2$.
Therefore
\begin{equation}\label{eq:anisotropic}
p\mid Q(x,y)\quad\Longleftrightarrow\quad p\mid x\text{ and }p\mid y.
\end{equation}
We first establish the unfiltered Hecke estimate. Section~\ref{sec:odd-ramified}
will impose oddness and derive the local bounds used in the sieve.

\subsection{The finite quotient and its Hecke operator}
Fix an odd $p\in\PP$, so $p\ge5$, and first suppose $p\nmid n$.
Reduction gives $\OO/p\OO\simeq M_2(\F_p)$, with norm equal to determinant.
Write
\[
G_p=\PGL_2(\F_p),\quad \Gamma_p=\rho_p(\OO^\times),\quad
Y_p=\Gamma_p\backslash G_p.
\]
Only $\pm1$ among the units reduce to scalars, so
$|\Gamma_p|=6$ and $|Y_p|=p(p^2-1)/6$. Equip $\C[Y_p]$ with counting
measure and define
\begin{equation}\label{eq:hecke}
(T_nf)(\Gamma_p g)=\frac1{12}\sum_{\alpha\in\mathcal Q_n}
f(\Gamma_p\rho_p(\alpha)g).
\end{equation}
Right multiplication by a unit permutes $\mathcal Q_n$, proving that this
is well-defined. It commutes with the right $G_p$ action. Conjugation
preserves $\mathcal Q_n$ and is projective inversion, so $T_n$ is
self-adjoint. Also $T_n1=\sigma(n)1$.

The determinant character
$\chi_p(\Gamma_p g)=(\det(g)/p)$ is well-defined, because the units have
norm one and scalar matrices have square determinant. Directly,
$T_n\chi_p=(n/p)\sigma(n)\chi_p$.

\subsection{The spectral bound}
This is the main imported ingredient. It turns the large supply of norm
representations into control of their distribution modulo a prime. The
normalization matters, so we identify the concrete operator before using
Jacquet--Langlands and the weight-two Ramanujan bound.

There are three things to check: the finite operator must be the stated
Hecke correspondence, the two large character components must be removed,
and the remaining spectrum must transfer to weight two. The proof below
checks these points in that order. Readers following the arithmetic first
can use the bound and continue with the normalizer average in the next
subsection.

\begin{lemma}\label{lem:spectral}
On $\{1,\chi_p\}^\perp\subset\C[Y_p]$, the operator $T_n$ has norm at most
$\tau(n)\sqrt n$, provided $(n,3p)=1$.
\end{lemma}
\begin{proof}
Class number one identifies the adelic quotient at the reduction kernel
with $\Gamma_p\backslash G_p$.  Indeed an adelic ideal principalizes;
and a rational projective element in all local order-unit groups can
be rescaled by a rational scalar to an order unit, since the local
valuations of its norm are even.  Its invariant measure is counting
measure up to normalization.

The all-integral-ideal correspondence divides the principal shell by
twelve units.  Principalization acts by $\rho_p(\alpha^{-1})$;
projectively this equals $\rho_p(\bar\alpha)$, and conjugation
permutes the shell.  Thus its operator is exactly \eqref{eq:hecke},
including the prime-power normalization.  Every global norm character
restricts to a character of $G_p$.  A character of
$\mathrm{GL}_2(\F_p)$ kills elementary unipotents and hence factors
through determinant; descent to $\mathrm{PGL}_2$ leaves only the
trivial and quadratic determinant characters.  Both have been
identified above.  Removing them, Lemma~\ref{lem:hecke-transfer}
therefore applies for $(n,3p)=1$ and gives the asserted bound.
\end{proof}

\subsection{Torus-normalizer averaging}
The image $C_p=(R/pR)^\times/\F_p^\times$ is a nonsplit torus. Its
normalizer is $K_p=C_p\cup C_p\rho_p(j)$, of order $2(p+1)$.
All twelve order units lie in $R^\times\cup R^\times j$, so
$\Gamma_p\subset K_p$. Put
\[
O_p=\Gamma_p\backslash K_p\subset Y_p,\qquad f_p=1_{O_p},\qquad
\mu_p=\frac{|O_p|}{|Y_p|}=\frac2{p(p-1)}.
\]
An invertible quaternion has one of its two Eisenstein components zero
modulo $p$ exactly when its projective reduction lies in $K_p$.
Let $b_p(n)$ count the unfiltered representations spoiled by $p$. Consequently
\[
b_p(n)=12(T_nf_p)(\Gamma_p).
\]
Both $f_p$ and $T_nf_p$ are right-$K_p$ invariant. Therefore
\[
(T_nf_p)(\Gamma_p)=\frac{\langle f_p,T_nf_p\rangle}{|O_p|}.
\]
The norm map $\F_{p^2}^\times\to\F_p^\times$ is onto, so $O_p$ has
equally many determinant-square and determinant-nonsquare points.
Thus $f_p-\mu_p1$ is orthogonal to both $1$ and $\chi_p$, and has
squared norm $|O_p|(1-\mu_p)$. The spectral lemma gives
\begin{equation}\label{eq:local}
\frac{b_p(n)}{r(n)}
\le\frac2{p(p-1)}+
\frac{\tau(n)\sqrt n}{\sigma(n)}
\left(1-\frac2{p(p-1)}\right).
\end{equation}
